\begin{figure}[ht]
\centering
\begin{tikzpicture}[>=Stealth, font=\small, node distance=2.5cm]
  \node[draw, rounded corners, fill=blue!8, minimum width=3.6cm, minimum height=1cm] (x1) {$X_0(N)$};
  \node[draw, rounded corners, fill=green!8, right=of x1, minimum width=4.4cm, minimum height=1cm] (z) {$Z_p=X_0(N;p)$};
  \node[draw, rounded corners, fill=blue!8, right=of z, minimum width=3.6cm, minimum height=1cm] (x2) {$X_0(N)$};

  \draw[->, thick] (z) -- node[above] {$\pi_1(E,C,D)=(E,C)$} (x1);
  \draw[->, thick] (z) -- node[above] {$\pi_2(E,C,D)=\bigl(E/D,(C+D)/D\bigr)$} (x2);

  \node[draw, rounded corners, fill=orange!10, below=1.8cm of z, minimum width=7.0cm, minimum height=0.95cm] (hecke)
    {$T_p=(\pi_2)_*\pi_1^*$，每个点有 $p+1$ 个阶 $p$ 子群};
  \draw[->, thick] (x1.south) -- ++(0,-0.7) -| (hecke.west);
  \draw[->, thick] (x2.south) -- ++(0,-0.7) -| (hecke.east);
\end{tikzpicture}
\caption{Hecke 算子不是凭空写下的矩阵，而是模曲线上的代数对应：先忘掉阶 $p$ 子群，再取商，最后在除子类与 Jacobian 上得到自同态。}
\label{fig:ch07-hecke-correspondences}
\end{figure}
